\section{Centera -- Content Addressed Storage}

Centera is a commercial quality object store that combines a
sophisticated software ``operating environment'' (OE) with low cost,
highly reliable hardware. The hardware is comprised of a set of discrete
nodes coupled together in such a manner that disruption in any one node
does not affect the operation of other nodes, yet appears to the
application (both physically and logically) as a single unit, or
cluster. This architecture, termed RAIN (redundant array of independent
nodes), allows for incremental scaling, from terabytes to petabytes,
without the need to disrupt system operation. Further, since the system
can continue to operate in the presence of failed nodes, service is
reduced to a regularly scheduled event where failed nodes are simply
replaced with new.

\subsection{Benefits of an Intrinsic Operating
Environment}

While the advancements in hardware that make such an architecture not
only possible but affordable are truly impressive, the majority of
features provided by Centera come from the operating environment,
CentraStar. CentraStar presents a flat object store to the application.
Applications submit data and are returned a global handle that uniquely
refers to the data stored. The handle returned is a function of the data
submitted. Upon submission, a 128-bit hash (called the Content Address,
or CA) is calculated at the application server. This value is then
recalculated upon arrival of the data to the CentraStar server. If the
values match the data is stored, otherwise an error is returned
indicating a corruption occurred during network transport. Thus the CA
serves the dual role of unique object handle and data checksum. The
checksum function is important as it serves as the means for integrity
checks throughout the life of the data. It's one thing for an
application to store mortgage application scans, it's of greater
importance that the data retrieved 30 years later can be validated to be
an accurate bit-for-bit representation of the data originally stored.

Data stored on Centera is automatically mirrored to multiple nodes. Thus
the failure of a node does not result in lost data. Rather, in the event
of a node failure all data that resided on that node is automatically
regenerated by CentraStar to assure there are always at least two copies
of all data objects. All of this occurs without any action by the
application or IT Ops. This model extends equally well when replicating
data between Centera clusters. Data stored on one Centera will
automatically be replicated to backup, again accomplished without human
intervention.

All of these actions are able to occur without human intervention or
setup due to the intrinsic intelligence that CentraStar brings to the
Centera offering.

\subsection{The Next Step Forward}

The combination of opaque object store and unique data handle (CA)
provide the needed next step forward in data abstraction. With Centera,
applications are successfully abstracted from the detail of data storage
and co-ownership of the data management function. IT Operations is
likewise abstracted from traditional storage details such as RAID types
to choose, LUNs to bind, or file systems to create \cite{ref8}.

This model of storage is named ``Content Addressed Storage'' or CAS. CAS
significantly extends the benefits of the object store contrasted to an
HFS in a number of significant areas.

\begin{enumerate}
\def\labelenumi{\arabic{enumi}.}
\item
  CAS creates a self-verifying namespace. With an HFS no such function
  is provided.
\item
  The object handle (Content Address) is guaranteed to return the data
  originally stored in the precise form that it was originally stored,
  irrespective of changes to the Centera configuration. HFS provides no
  such guarantee.
\item
  Unlike CAS, the HFS pathname holds no intrinsic semantic meaning.
\item
  With CAS the namespace, as seen by the application, is guaranteed not
  to change. With HFS namespace changes are not only possible, but
  frequent.
\end{enumerate}
